\section{Add One to a Number Represented by a Linked List}
\label{sec:ll-add-one}

\problemstatement{A non-negative integer is stored as a linked list, most
significant digit first (head). Add one to the number and return the head of
the resulting list.}

\begin{patternbox}
The carry propagates from the \emph{least} significant digit, but the list
is stored most-significant first -- so either \textbf{reverse, add, reverse
back}, or use recursion to reach the last node and carry upward on the
return trip. Both avoid needing backward pointers.
\end{patternbox}

\subsection*{Carry from the back via recursion}

Recurse to the tail; the deepest call adds one and returns any carry. Each
frame adds the incoming carry to its digit, updates the digit modulo $10$,
and returns the new carry ($1$ if the digit was $9$, else $0$). If the head
still produces a carry after processing, prepend a new leading $1$ node.
This treats the recursion stack as the ``reversed'' traversal, so no
explicit reversal is required.

\begin{ideabox}[The Return Trip Is the Least-Significant-First Pass]
Recursion reaches the last (least significant) digit first on the way down,
then processes digits back-to-front on the way up -- exactly the order
addition needs. A leftover carry at the very top means a new most-
significant digit.
\end{ideabox}

\subsection*{Algorithm}

\begin{enumerate}
  \item Recurse to the last node; add $1$ there, return the carry.
  \item Each frame: \texttt{sum = digit + carry}; set digit to
        \texttt{sum \% 10}; return \texttt{sum / 10}.
  \item If the top-level carry is $1$, prepend a new node with value $1$.
\end{enumerate}

\subsection*{Dry run}

\dryrun{List $1\to2\to9$ (number $129$).}

\begin{itemize}
  \item Last digit $9+1=10$: digit $0$, carry $1$. Next $2+1=3$, carry $0$.
  \item Head $1$ unchanged. Result $1\to3\to0$ ($130$).
\end{itemize}

\subsection*{C++ implementation}

\begin{lstlisting}
#include <bits/stdc++.h>
using namespace std;

struct ListNode {
    int val; ListNode* next;
    ListNode(int v) : val(v), next(nullptr) {}
};

ListNode* build(const vector<int>& a) {
    ListNode dummy(0), *t = &dummy;
    for (int x : a) { t->next = new ListNode(x); t = t->next; }
    return dummy.next;
}

int addHelper(ListNode* node) {
    if (!node) return 1;                          // seed the +1 at the end
    int carry = addHelper(node->next);
    int sum = node->val + carry;
    node->val = sum % 10;
    return sum / 10;
}

ListNode* addOne(ListNode* head) {
    int carry = addHelper(head);
    if (carry) {                                  // new leading digit
        ListNode* node = new ListNode(carry);
        node->next = head;
        return node;
    }
    return head;
}

void print(ListNode* h) { for (; h; h = h->next) cout << h->val; cout << "\n"; }

int main() {
    print(addOne(build({1, 2, 9})));              // 130
    print(addOne(build({9, 9, 9})));              // 1000
    return 0;
}
\end{lstlisting}

\begin{complexitybox}
\textbf{Time Complexity.} $O(n)$. \textbf{Space Complexity.} $O(n)$
recursion stack (or $O(1)$ with the reverse--add--reverse approach).
\end{complexitybox}

\begin{takeawaybox}
Arithmetic on a most-significant-first list needs a back-to-front pass,
supplied for free by recursion's return trip (or by reversing). A leftover
carry at the head becomes a new leading digit -- the one edge case to
remember.
\end{takeawaybox}
